\subsection{Beyond Classification}
Although the majority of the work on fairness in machine learning focuses on batch classification, batch classification is only one aspect of how machine learning is used. Much of machine learning --- e.g. online learning, bandit learning, and reinforcement learning --- focuses on \emph{dynamic} settings in which the actions of the algorithm feed back into the data it observes. These dynamic settings capture many problems for which fairness is a concern. For example, lending, criminal recidivism prediction, and sequential drug trials all are so-called \emph{bandit} learning problems, in which the algorithm cannot observe data corresponding to counterfactuals. We cannot see whether someone not granted a loan \emph{would have} paid it back. We cannot see whether an inmate not released on parole \emph{would have} gone on to commit another crime. We cannot see how a patient \emph{would have} responded to a different drug.

The theory of learning in bandit settings is well understood, and it is characterized by a need to trade off \emph{exploration} with \emph{exploitation}. Rather than always making a myopically optimal decision, when counter-factuals cannot be observed, it is necessary for algorithms to sometimes take actions that appear to be sub-optimal so as to gather more data. But in settings in which decisions correspond to individuals, this means sacrificing the well-being of a particular person for the potential benefit of future individuals. This can sometimes be unethical, and a source of unfairness \cite{explore}. Several recent papers explore this issue. For example, \cite{BBK17} and \cite{KMRWW18} give conditions under which linear learners need not explore at all in bandit settings, thereby allowing for best-effort service to each arriving individual, obviating the tension between ethical treatment of individuals and learning. \cite{RSVW18} show that the costs associated with exploration can be unfairly bourn by a structured sub-population, and that counter-intuitively, those costs can actually increase when they are included with a majority population, even though more data increases the rate of learning overall. However, these results are all preliminary: they are restricted to settings in which the learner is learning a linear policy, and \emph{the data really is governed by a linear model}. While illustrative, more work is needed to understand real-world learning in online settings, and the ethics of exploration.

There is also some work on fairness in machine learning in other settings --- for example, ranking \cite{YS17,CSV17}, selection \cite{selection,KR18}, personalization \cite{celis2017fair}, bandit learning \cite{infinitebandit,calibratedbandit}, human-classifier hybrid decision systems \cite{defer}, and reinforcement learning \cite{JJKMR17,DTB17}. But outside of classification, the literature is relatively sparse. This should be rectified, because there are interesting and important fairness issues that arise in other settings --- especially when there are combinatorial constraints on the set of individuals that can be selected for a task, or when there is a temporal aspect to learning.
